\section{gepaLLMAsJudge}

The \texttt{gepaLLMAsJudge} workflow evaluates whether prompt optimization
improves a local judge for
the seven adapted Schoenfeld episode labels: Read, Analyze, Plan, Implement,
Explore, Verify, and Monitor.  The model was the Q4\_K\_XL quantization of
Qwen3.6-27B served by \texttt{llama.cpp} on one RTX 3090.  All splits preserve
complete responses, preventing neighboring units from crossing partitions.

\paragraph{Preliminary single-split runs.}
With seed 42, the 38 responses were divided into 26 training, 6 validation,
and 6 test documents (2,382/407/336 units).  These first runs classified the
current unit without neighboring context and optimized per-unit exact match.

\begin{table}[t]
\centering
\caption{Preliminary \texttt{gepaLLMAsJudge} agreement.  The composite rescales Cohen's
$\kappa$ and Kendall's $\tau_b$ from $[-1,1]$ to $[0,1]$ and averages them.
Coverage was 100\% in every row.  \modelmetrics{Qwen3.6-27B
\texttt{Q4\_K\_XL}, used as the episode-label judge}{exact-match accuracy,
Cohen's $\kappa$, Kendall's $\tau_b$, their rescaled composite, and valid-output
coverage}}
\label{tab:exp0a-judge}
\begin{tabular}{lrrrr}
\toprule
Evaluation & Accuracy & $\kappa$ & $\tau_b$ & Composite \\
\midrule
Seed validation      & 0.6388 & 0.5581 & 0.4810 & 0.7598 \\
Optimized validation & 0.6855 & 0.6074 & 0.5847 & 0.7980 \\
Optimized test       & 0.6607 & 0.5710 & 0.6066 & 0.7944 \\
Few-shot validation  & 0.6806 & 0.6078 & 0.5830 & 0.7977 \\
Few-shot test        & \textbf{0.6994} & \textbf{0.6178} &
\textbf{0.6667} & \textbf{0.8211} \\
\bottomrule
\end{tabular}
\end{table}

The matched few-shot run used the same seed, document split, model, and GEPA
budget, adding seven gold training examples to the instruction.  It performs
better than the optimized base condition on the held-out test set: accuracy
rises from 66.07\% to 69.94\% (+3.87 percentage points), Cohen's $\kappa$ from
0.5710 to 0.6178, Kendall's $\tau_b$ from 0.6066 to 0.6667, and composite
agreement from 0.7944 to 0.8211.  GEPA itself does not improve the few-shot
validation result: the seed and selected optimized prompts both score 68.06\%
accuracy.  The current comparison therefore favors few-shot prompting for
held-out generalization, whereas the measurable GEPA optimization gain is
confined to the base-prompt condition.

\paragraph{Context-aware nested cross-validation.}
The revised protocol excludes the same six responses from model selection and
runs five outer folds over the remaining 32 responses.  Within each fold,
five responses form the inner validation set and all others outside the outer
fold form the GEPA training set.  Each request includes the problem statement,
previous unit, current unit, and next unit.  The seed instruction includes 21
audited synthetic contrastive demonstrations.  GEPA receives inverse-frequency
weighted exact match, whose split mean is balanced accuracy; the safety
selector compares either balanced accuracy or macro-F1 and rejects any
candidate whose class recall falls by more than the configured threshold.

\paragraph{Informative same-model reward.}
The current implementation can replace the static GEPA reward with an
LLM-as-a-judge call to the same Qwen3.6-27B model.  For each candidate label,
the judge returns a short free-text explanation plus four diagnostic scores
from 1 to 5: gold agreement, functional fit, contextual coherence, and boundary
precision.  Their normalized weighted mean is used only as GEPA's optimization
signal, with respective weights 0.40, 0.25, 0.15, and 0.20.  Gold-label
accuracy, balanced accuracy, macro-F1, and per-class recall remain the
independent selection and reporting criteria.  The judge runs deterministically
with model reasoning disabled.  Every judge call is persisted as JSONL with
its component scores, normalized score, and free-text explanation, permitting
a direct audit of feedback quality.  Tables~\ref{tab:exp0a-cv} and
\ref{tab:exp0a-classes} remain static-reward cross-validation baselines; the
completed LLM-judge final fit is reported separately below because its reused
validation set is a prompt-development set rather than an outer-fold estimate.

\begin{table}[t]
\centering
\caption{Aggregate selected-prompt outer-fold results over 2,789 units.
``Balanced light'' is the earlier 0.10 recall-gate run; the macro-F1 runs use a
0.05 gate.  The six excluded responses were not evaluated.
\modelmetrics{Qwen3.6-27B \texttt{Q4\_K\_XL}, used as the episode-label
judge}{exact-match accuracy, balanced accuracy, macro-F1, and the rescaled
$\kappa$/$\tau_b$ composite}}
\label{tab:exp0a-cv}
\begin{tabular}{lrrrr}
\toprule
Selection / budget & Accuracy & Balanced acc. & Macro-F1 & Composite \\
\midrule
Balanced / light & 0.6701 & 0.6753 & 0.6480 & 0.8031 \\
Macro-F1 / light & \textbf{0.6734} & \textbf{0.6796} &
\textbf{0.6505} & \textbf{0.8063} \\
Macro-F1 / heavy & 0.6669 & 0.6737 & 0.6437 & 0.8030 \\
\bottomrule
\end{tabular}
\end{table}

Macro-F1 selection with the light budget is the strongest cross-validated
configuration, but its advantage over balanced-accuracy selection is small:
+0.32 accuracy points, +0.43 balanced-accuracy points, and +0.24 macro-F1
points.  At the response level it wins on 15 documents, loses on 13, and ties
on 4, so the observed gain is not stable across documents.  The light run
selects a GEPA-optimized prompt in three of five folds; the recall gate retains
the seed prompt in the other two because Plan or Monitor recall falls too far.

Increasing the budget does not help.  The heavy run selects an optimized
prompt in only two folds and is 0.65 accuracy points and 0.68 macro-F1 points
below the light run.  Paired predictions show 38 units correct only under
light versus 20 correct only under heavy.  This difference is concentrated in
fold 5, where light discovers and selects a strong candidate whereas heavy
retains the seed after its candidate violates the Analyze recall gate.
Independent heavy search is therefore not a superset of light search and can
discard a better prompt found by a smaller run.

\begin{table}[t]
\centering
\caption{Per-class outer-fold performance for the best current configuration,
macro-F1 selection with the light GEPA budget.  \modelmetrics{Qwen3.6-27B
\texttt{Q4\_K\_XL}, used as the episode-label judge}{class support, precision,
recall, and F1}}
\label{tab:exp0a-classes}
\begin{tabular}{lrrrr}
\toprule
Class & Support & Precision & Recall & F1 \\
\midrule
Read      & 318 & 0.798 & 0.770 & 0.784 \\
Analyze   & 685 & 0.757 & 0.632 & 0.689 \\
Plan      & 205 & 0.524 & 0.590 & 0.555 \\
Implement & 650 & 0.880 & 0.722 & 0.793 \\
Explore   & 193 & 0.679 & 0.492 & 0.571 \\
Verify    & 538 & 0.669 & 0.606 & 0.636 \\
Monitor   & 200 & 0.364 & 0.945 & 0.526 \\
\bottomrule
\end{tabular}
\end{table}

The revised objective improves overall balance only marginally and does not
resolve Monitor overprediction.  The model emits Monitor 519 times for 200
gold Monitor units; the leading errors are Verify$\rightarrow$Monitor (127),
Plan$\rightarrow$Monitor (65), and Explore$\rightarrow$Monitor (42).  The
annotation audit flags 180/3,125 units (5.76\%), including 145 multiline units
and 49 structural/control markers such as \texttt{</think>} and ``Final
Answer.''  These artifacts plausibly inflate Monitor recall while depressing
precision and should be reviewed before further optimization.

All runs retain 100\% final-output coverage and complete without OOM failures,
but the light and heavy jobs log 31 and 28 truncated LM responses during
optimization.  Moreover, the outer folds currently evaluate only the
safety-selected prompt, not both seed and optimized prompts, so they do not
directly estimate GEPA's paired outer-fold gain.  Finally, the six responses
excluded from current CV were inspected in the preliminary experiments; they
are locked from the current code path but no longer constitute a pristine
research holdout.  The judge should therefore not replace validated annotation
until structural artifacts are adjudicated, decoding is constrained, paired
outer evaluation is added, and a genuinely unseen external test set is
designated.

\paragraph{Completed LLM-judge final fit.}
The August 11, 2026 run uses seed 42, the 21-example context-aware few-shot
prompt, and the same-model informative reward on the original response-grouped
26/6/6 split (2,382/407/336 units).  Seed and optimized prompts both return a
valid label for all 407 validation units.  Table~\ref{tab:exp0a-llmjudge}
reports exact gold-label metrics rather than the LLM judge's diagnostic score.

\begin{table*}[t]
\centering
\caption{Context-aware LLM-judge final-fit results on the 407-unit
prompt-development validation set.  Coverage is 100\% in both rows.  The locked test
is not evaluated.  \modelmetrics{Qwen3.6-27B \texttt{Q4\_K\_XL}, used both
for episode-label prediction and the same-model GEPA reward}{exact-match
accuracy, balanced accuracy, macro-F1, Cohen's $\kappa$, Kendall's $\tau_b$,
absolute change, and valid-output coverage}}
\label{tab:exp0a-llmjudge}
\begin{tabular}{lrrrrr}
\toprule
Prompt & Accuracy & Balanced acc. & Macro-F1 & $\kappa$ & $\tau_b$ \\
\midrule
21-example seed & 0.6683 & 0.6765 & 0.6132 & 0.5988 & 0.5904 \\
GEPA optimized  & \textbf{0.7543} & \textbf{0.7769} & \textbf{0.7270} &
\textbf{0.6968} & \textbf{0.6832} \\
\midrule
Absolute change & +0.0860 & +0.1003 & +0.1138 & +0.0980 & +0.0929 \\
\bottomrule
\end{tabular}
\end{table*}

The optimized prompt fixes 49 seed errors and regresses on 14 previously
correct units, a net gain of 35.  All six validation responses improve, with
response-level accuracy gains from 2.0 to 13.8 percentage points.  Its primary
effect is to curb the seed prompt's overuse of Monitor: Monitor predictions
fall from 74 to 43, including nine corrections from Monitor to true Plan and
15 from Monitor to true Verify.  Recall rises most for Plan (+28.13 points),
Explore (+23.53), Verify (+14.29), and Analyze (+8.20).  Monitor recall falls
slightly from 1.000 to 0.944, but precision rises from 0.243 to 0.395.  The
selected prompt still has weak Monitor precision and Analyze recall, and the
remaining errors concentrate at the Analyze/Implement/Verify boundaries and
Verify$\rightarrow$Monitor.  Implement is the only class whose F1 decreases,
from 0.839 to 0.825, because its additional predictions include more false
positives.

\begin{table}[t]
\centering
\caption{Per-class recall and optimized F1 for the completed LLM-judge final
fit.  \modelmetrics{Qwen3.6-27B \texttt{Q4\_K\_XL}, used both for
episode-label prediction and the same-model GEPA reward}{class support, seed
and optimized recall, and optimized per-class F1}}
\label{tab:exp0a-llmjudge-classes}
\begin{tabular}{lrrrr}
\toprule
Class & Support & Seed recall & Opt. recall & Opt. F1 \\
\midrule
Read      &  45 & 0.778 & 0.756 & 0.791 \\
Analyze   & 122 & 0.598 & 0.680 & 0.744 \\
Plan      &  32 & 0.531 & 0.813 & 0.722 \\
Implement & 103 & 0.786 & 0.825 & 0.825 \\
Explore   &  17 & 0.471 & 0.706 & 0.686 \\
Verify    &  70 & 0.571 & 0.714 & 0.763 \\
Monitor   &  18 & 1.000 & 0.944 & 0.557 \\
\bottomrule
\end{tabular}
\end{table}

The internal same-model reward rises from 0.7264 for the seed to 0.8196 for
the selected candidate.  The optimization trace contains 145 candidates, 144
full validation evaluations, and 60,767 persisted judge calls, of which 148
are judge-output errors.  A text audit finds 37 validation units for which the
judge explicitly disputes the corpus gold label in at least one evaluation,
showing that the reward itself detects annotation-sensitive boundaries rather
than consistently treating gold as semantically authoritative.  The structural
audit flags 29 validation units; their accuracy remains 72.4\% under both
prompts, so the gain comes entirely from otherwise unflagged units.

GEPA expands the instruction from 21 to 40 demonstrations and from 7,648 to
32,717 characters.  The extra rules and examples target observed corpus edge
cases, raising inference cost and the risk of learning annotation conventions
rather than a portable episode taxonomy.  More importantly, GEPA evaluates
145 candidates against the same six-response validation set used for final
selection.  The nominal 407 units are also correlated within responses, and
only one optimization seed is available.  The result files explicitly record
\texttt{locked\_test\_evaluated: false} and \texttt{test: null}.  The observed
gain is therefore strong evidence of successful prompt optimization on the
development distribution, but not a held-out generalization estimate.  The
prompt should be frozen and compressed or ablated before a one-time test on
genuinely unseen responses.

\paragraph{Final Qwen3.8 model-alignment pass.}
The August 27--28, 2026 alignment run uses Qwen3.8-27B
\texttt{UD-Q4\_K\_XL}, seed 42, the heavy GEPA budget, and medium reasoning
effort for the same-model judge.  It keeps the same response-grouped
2,382/407/336 train/validation/locked-test split.  Rather than drawing new
few-shot examples, it imports the previous 40-example Qwen3.6 optimized prompt
as its external seed.  Gold-label agreement supplies 40\% of the internal
judge reward; functional fit, contextual coherence, and boundary precision
supply the remaining 60\%.  Final selection uses exact balanced accuracy with
a maximum allowed per-class recall drop of 0.10.

\begin{table*}[t]
\centering
\caption{Final Qwen3.8 model-alignment results on the 407-unit development
validation set.  Both prompts have 100\% valid-output coverage; the 336-unit
locked test is not evaluated.  \modelmetrics{Qwen3.8-27B
\texttt{UD-Q4\_K\_XL}, used for both classification and the same-model GEPA
reward}{exact-match accuracy, balanced accuracy, macro-F1, Cohen's $\kappa$,
and Kendall's $\tau_b$}}
\label{tab:exp0a-qwen38-final}
\begin{tabular}{lrrrrr}
\toprule
Prompt & Accuracy & Balanced acc. & Macro-F1 & $\kappa$ & $\tau_b$ \\
\midrule
Imported 40-example seed & \textbf{0.7174} & \textbf{0.7400} &
\textbf{0.6938} & \textbf{0.6523} & \textbf{0.6122} \\
GEPA optimized & 0.7125 & 0.7288 & 0.6743 & 0.6475 & 0.6009 \\
\midrule
Absolute change & $-0.0049$ & $-0.0112$ & $-0.0196$ & $-0.0049$ & $-0.0113$ \\
\bottomrule
\end{tabular}
\end{table*}

\begin{table}[t]
\centering
\caption{Per-class recall in the final Qwen3.8 alignment.  The selection gate
rejects any candidate whose recall drops by more than 0.10 for any class.
\modelmetrics{Qwen3.8-27B \texttt{UD-Q4\_K\_XL}}{class support and seed versus
GEPA-optimized recall}}
\label{tab:exp0a-qwen38-classes}
\begin{tabular}{lrrr}
\toprule
Class & Support & Seed recall & Opt. recall \\
\midrule
Read      &  45 & 0.778 & 0.756 \\
Analyze   & 122 & 0.549 & 0.533 \\
Plan      &  32 & 0.750 & 0.563 \\
Implement & 103 & 0.874 & 0.874 \\
Explore   &  17 & 0.529 & 0.647 \\
Verify    &  70 & 0.700 & 0.786 \\
Monitor   &  18 & 1.000 & 0.944 \\
\bottomrule
\end{tabular}
\end{table}

The optimized candidate changes 49 labels: it fixes 17 seed errors, breaks 19
seed-correct cases, and changes 13 errors to other errors.  Its largest failure
is Plan recall, which falls by 0.188 and therefore breaches the recall gate;
the selected prompt is byte-identical to the imported seed.  An unclustered
exact McNemar test on the 17 versus 19 discordant correctness outcomes gives
$p=0.868$, offering no evidence of an aggregate improvement, although this
unit-level calculation understates response clustering.

This negative exact-label result contrasts with the internal GEPA objective,
whose best candidate reward rises from 0.7684 to 0.7762.  Because only 40\% of
that reward is gold agreement, the discrepancy is consistent with optimizing
the judge's broader functional criteria without improving corpus-label
classification.  The optimized instruction is shorter (2,591 versus 5,201
words) but internally inconsistent: an early rule treats ``Let me compute
that'' as Monitor while a later rule treats ``Let me calculate'' as Plan.
Plan-to-Monitor errors rise from 6 to 10, consistent with that boundary
confusion.  The best internal candidate appears by iteration 3 and remains the
best through iteration 19, so the remaining heavy search adds roughly 14.5
hours without a better candidate.  The appropriate conclusion is that the
final alignment completed correctly and the safety gate worked: GEPA found no
safe improvement over the existing seed.  The complete artifact is stored at
\path{results/gepaLLMAsJudge/qwen3.8-27b-medium-final-s42-v3}.
